\documentclass[11pt,a4paper]{article}
\usepackage[margin=25mm]{geometry}
\usepackage{amsmath,amssymb,amsthm}
\newtheorem{theorem}{Theorem}
\title{Regular polytopes for four alternating groups}
\author{Human Compression project: L071 recovery}
\date{Working statement; frozen scientific review pending}
\begin{document}\maketitle
The complete retained proof of enumeration by stabilizers of initial
strings supplies the common theory. In this note isomorphism preserves
the order of the distinguished generators. A marking is therefore
counted separately from its dual unless an automorphism reverses it.
This is the convention of Leemans--Mixer, \emph{Algorithms for
classifying regular polytopes with a fixed automorphism group},
Contributions to Discrete Mathematics 7(2) (2012), 105--118,
Section~6 and Table~2, physical page~9.

\begin{theorem}
The numbers of regular polytopes with the indicated automorphism
groups, up to isomorphism, are as follows. Every unlisted rank has
zero polytopes.
\[
\begin{array}{c|rrrr}
 &\text{rank }3&\text{rank }4&\text{rank }5&\text{rank }6\\\hline
 A_{11}&122&0&0&4\\
 A_{12}&375&173&43&0\\
 A_{13}&3021&202&50&20\\
 A_{14}&8531&255&90&18
\end{array}
\]
In particular the missing ranks four and five for $A_{11}$ are genuine
gaps in its rank spectrum.
\end{theorem}
\begin{proof}
For $11\le n\le14$ take $G=A_n$ in its natural permutation action.
We use the classical automorphism theorem
$\operatorname{Aut}(A_n)=S_n$ acting by conjugation for $n\ge7$.
Thus the conjugation action of $S_n$ supplies precisely the abstract
isomorphisms between markings on $G$. This standard theorem is a
retained input, not a conclusion of the finite search.

The program forms the complete set of nonidentity involutions of
$A_n$ from its conjugacy classes. At an accepted prefix it chooses
one representative of each orbit of its pointwise stabilizer in
$S_n$ on the eligible new involutions, and applies the exact end
intersection test. Its cached commuting sets and subgroup tests are
the retained common algorithm. The implementation also records
prefixes generating other transitive subgroups, and continues extending
those prefixes. We retain only the leaves generating $A_n$ for the
present claim. These extra records do not merge or duplicate any
prefix orbit, and the algorithm stops extending a prefix that already
generates $A_n$.

One implementation detail weakens, without invalidating, the common
pruning test. A branch is discarded only when the group generated by
the old prefix and every allowed commuting involution is intransitive.
Every string extension lies in this group. An intransitive group cannot
contain $A_n$, so no generating $A_n$ marking is lost. The test may
retain extra branches, which does not affect completeness or uniqueness.
Every appended mark is outside the old generated subgroup. Hence the
subgroup-index argument bounds the rank by
$\lfloor\log_2|A_n|\rfloor$, giving the respective complete search
limits $24,27,31,35$. The complete runs retain all literal tuples and
all zero rows through these limits.

The program records respectively $136,3935,3293,9562$ transitive
markings. Exact generation tests select $126,591,3293,8894$ of these
as generating $A_{11},A_{12},A_{13},A_{14}$. Grouping those selected
tuples by length gives the displayed table. In addition to the
recursive enumeration test, a separate execution checks every pair of
subsets of the generators of every recorded transitive tuple against
the full intersection property. Those executions also check involution
orders, string commutations and transitivity, and repeat the exact
generation tests. All pass. These additional checks verify the recorded
objects; completeness rests on the retained orbit enumeration proof and
the completed full-rank executions.

Ranks one and two cannot occur. A group generated by one involution
has order at most two. A group generated by two involutions has a cyclic
subgroup of index at most two, generated by their product. Neither can
be the nonabelian simple group $A_n$. This covers the ranks below the
enumeration's recording threshold and completes the proof.
\end{proof}

\paragraph{Attribution, finite premises and remaining source scope.}
The table recovers exactly the original Table~2; it is not a new
classification claim. The stabilizer method is the established method
described in the source and attributed there to Hartley--Hulpke.
Our executable proof makes the extension criterion, pruning and rank
coverage explicit and retains the full finite output. All group
operations are exact GAP operations and remain software premises.
There is no claim of Lean verification.

The source's historical conjecture in Remark~6.1 is separate from these
four finite classifications. Its later resolution must not be credited
to this computation. Likewise the source's compatible-action theorem,
algorithm comparisons, smaller group classifications, sporadic
verification and full $\mathrm{Co}_3$ tables are separate obligations.
This component alone is not a recovery of the whole source and supplies
no whole-original or compression-factor credit.
\end{document}
